要求积性函数 $f$ 的前缀和
\[
  S_f(n)=\sum_{i=1}^{n}f(i).
\]
设
\[
  f*g=h,
\]
即
\[
  h(n)=\sum_{d\mid n}f(d)g\!\left(\frac nd\right).
\]

对 $h(n)$ 求前缀和：
\[
  \begin{aligned}
  \sum_{i=1}^{n}h(i)
  &=\sum_{i=1}^{n}\sum_{d\mid i}g(d)f\!\left(\frac id\right)\\
  &=\sum_{d=1}^{n}g(d)
    \sum_{k=1}^{\lfloor n/d\rfloor}f(k)
    \qquad(i=dk)\\
  &=\sum_{d=1}^{n}g(d)S_f\!\left(\left\lfloor\frac nd\right\rfloor\right).
  \end{aligned}
\]

我们想求的是 $S_f(n)$，一般选择满足 $g(1)=1$ 的卷积，于是
\[
  \sum_{i=1}^{n}h(i)
  =S_f(n)+\sum_{d=2}^{n}g(d)
    S_f\!\left(\left\lfloor\frac nd\right\rfloor\right),
\]
从而
\[
  S_f(n)
  =\sum_{i=1}^{n}h(i)-\sum_{d=2}^{n}g(d)
    S_f\!\left(\left\lfloor\frac nd\right\rfloor\right).
\]

因此要找一个卷积 $f*g=h$，使得 $h$ 的前缀和非常容易计算。
递推的后一部分使用整除分块，再从较小的
$S_f(\lfloor n/d\rfloor)$ 递归反推出 $S_f(n)$。

通常先用线性筛预处理 $B=\lfloor n^{2/3}\rfloor$ 范围内的前缀和，
避免大量小值被重复计算；超过 $B$ 的状态使用哈希表记忆化。

\subsubsection{例 1：$\mu$ 的前缀和}

要求
\[
  S_\mu(n)=\sum_{i=1}^{n}\mu(i).
\]
由
\[
  \mu*\mathbf 1=\varepsilon
\]
代入通式可得
\[
  \begin{aligned}
  S_\mu(n)
  &=\sum_{i=1}^{n}\varepsilon(i)
    -\sum_{d=2}^{n}S_\mu\!\left(\left\lfloor\frac nd\right\rfloor\right)\\
  &=1-\sum_{d=2}^{n}S_\mu\!\left(\left\lfloor\frac nd\right\rfloor\right).
  \end{aligned}
\]

\subsubsection{例 2：$\varphi$ 的前缀和}

要求
\[
  S_\varphi(n)=\sum_{i=1}^{n}\varphi(i).
\]
由
\[
  \varphi*\mathbf 1=\operatorname{id}
\]
代入通式可得
\[
  \begin{aligned}
  S_\varphi(n)
  &=\sum_{i=1}^{n}\operatorname{id}(i)
    -\sum_{d=2}^{n}S_\varphi\!\left(\left\lfloor\frac nd\right\rfloor\right)\\
  &=\frac{n(n+1)}2
    -\sum_{d=2}^{n}S_\varphi\!\left(\left\lfloor\frac nd\right\rfloor\right).
  \end{aligned}
\]
这里卷积的第二个函数是 $\mathbf 1$，所以 $g(d)=1$，递推项前没有额外的 $d$。

\subsubsection{整除分块}

对固定的 $n$，$\lfloor n/d\rfloor$ 只有 $O(\sqrt n)$ 种不同取值。
若当前左端点为 $l$，令
\[
  q=\left\lfloor\frac nl\right\rfloor,
  \qquad
  r=\left\lfloor\frac nq\right\rfloor,
\]
则所有 $d\in[l,r]$ 都满足 $\lfloor n/d\rfloor=q$，可以一次合并整段贡献。

\begin{icpcwarning}[实现细节]
  $n(n+1)/2$ 和分块贡献可能在乘法阶段溢出，模板使用
  \texttt{\_\_int128} 保存中间结果，再转回 \texttt{ll}。
  构造时的预处理上界至少应为 $1$。
\end{icpcwarning}
